% !PRINTED 542
\weberpage{542}
%<<CONT
modulus $\mu$ when the integers $\xi, \xi'$ are congruent. From this
it follows that if $\xi$ runs through a complete system of residues
with respect to the modulus $\mu$, the same also holds of the product
$\alpha\xi$, whereby the theorem is proved:

1. \gesperrt{If $\alpha$ is an integer in $\Omega$, relatively prime
to the modulus $\mu$, and further $\gamma$ is an arbitrary integer in
$\Omega$, then the congruence:}
\weq{1}{\alpha\xi \equiv \gamma \ (\mathrm{mod}\ \mu)}
\gesperrt{is always solvable by an integer $\xi$, and also by only
one, if a complete system of residues with respect to the modulus
$\mu$ is prescribed for $\xi$.}

If in this $\mu$ is itself a number, which we denote by $\beta$, then
the quotient
\[
\frac{\alpha\xi - \gamma}{\beta} = -\eta
\]
is also an integer, and the foregoing theorem then takes the form:

2. \gesperrt{If $\alpha, \beta, \gamma$ are three integers in $\Omega$
and $\alpha, \beta$ are relatively prime, then one can determine two
other integers $\xi, \eta$ in $\Omega$ so that}
\weq{2}{\alpha\xi + \beta\eta = \gamma}
\gesperrt{holds. In particular, one can therefore also, for two
arbitrary relative primes $\alpha, \beta$, satisfy the equation}
\weq{3}{\alpha\xi + \beta\eta = 1}
\gesperrt{by integers $\xi, \eta$.}

This theorem can be carried over to a system of several numbers.

If the numbers $\alpha, \beta, \gamma, \ldots$ in $\mathfrak{o}$ have
no common divisor, then we can first determine numbers $\eta_1, \zeta_1,
\ldots$ in $\mathfrak{o}$ so that
\[
\beta_1 = \beta\eta_1 + \gamma\zeta_1 + \cdots
\]
becomes relatively prime to $\alpha$. For if $\pi_1, \pi_2,
\ldots$ are the distinct prime factors of $\alpha$, no one of which
can divide all of $\beta, \gamma, \ldots$, and if, say, $\beta$ is not
divisible by $\pi_1$, we have to take $\eta_1$ indivisible by $\pi_1$,
the remaining numbers $\zeta_1, \ldots$ divisible by $\pi_1$, and so
forth, and we obtain for each of the numbers $\eta_1, \zeta_1, \ldots$
the condition that it be divisible by some of the prime factors
%CONT>>
